% !TEX program = xelatex
\documentclass[10pt, letterpaper]{article}

\usepackage[
  top=0.50in,
  bottom=0.55in,
  left=0.60in,
  right=0.55in
]{geometry}
\usepackage{fontspec}
\usepackage{paracol}
\usepackage{tabularx}
\usepackage{array}
\usepackage{enumitem}
\usepackage{fontawesome5}
\usepackage[hidelinks]{hyperref}
\usepackage{needspace}
\usepackage{xcolor}

\defaultfontfeatures{Ligatures=TeX}
\IfFontExistsTF{Tinos}{
  \setmainfont{Tinos}
}{
  \IfFontExistsTF{Times New Roman}{
    \setmainfont{Times New Roman}
  }{
    \setmainfont{TeX Gyre Termes}
  }
}

\pagestyle{empty}
\setlength{\parindent}{0pt}
\setlength{\parskip}{0pt}
\setlength{\columnsep}{0.10in}
\setlength{\columnseprule}{0pt}
\urlstyle{same}
\emergencystretch=2em

\newcommand{\squarebullet}{\raisebox{0.25ex}{\rule{0.45ex}{0.45ex}}}
\setlist[itemize]{
  label=\squarebullet,
  leftmargin=1.15em,
  topsep=0.16em,
  itemsep=0.05em,
  parsep=0pt,
  partopsep=0pt
}

\newcommand{\sectionlabel}[1]{%
  {\fontsize{9.5pt}{11pt}\selectfont\bfseries\MakeUppercase{#1}}%
}

\newenvironment{cvsection}[1]{%
  \par\vspace{0.86\baselineskip}%
  \Needspace{3\baselineskip}%
  \setcolumnwidth{1.10in,\fill}%
  \begin{paracol}{2}%
  \raggedright\sectionlabel{#1}%
  \switchcolumn
}{%
  \end{paracol}%
}

\newcommand{\externalicon}[1]{%
  \href{#1}{\raisebox{0.12ex}{\normalfont\scriptsize\faExternalLink*}}%
}

\newcommand{\educationentry}[4]{%
  \noindent\begin{tabularx}{\linewidth}{@{}Xr@{}}
    \textbf{#1} & \textbf{#2} \\
    \textit{#3} & \textit{\small #4}
  \end{tabularx}%
}

\newcommand{\researchentry}[2]{%
  \noindent\begin{tabularx}{\linewidth}{@{}Xr@{}}
    \textbf{#1} & \textit{#2}
  \end{tabularx}%
}

\newcommand{\projectentry}[2]{%
  \noindent\begin{tabularx}{\linewidth}{@{}Xr@{}}
    \textbf{#1} & \textit{#2}
  \end{tabularx}%
}

\newcommand{\awardentry}[2]{%
  \noindent\begin{tabularx}{\linewidth}{@{}p{1.15em}Xr@{}}
    \squarebullet & #1 & \textit{#2}
  \end{tabularx}%
}

\begin{document}

\begin{center}
  {\fontsize{24pt}{27pt}\selectfont\bfseries Yibo Yin}

  \vspace{0.58em}
  Email: \href{mailto:yibo.yin@epfl.ch}{yibo.yin@epfl.ch}
  \quad\textbullet\quad
  Phone: \href{tel:+8613948812636}{(+86) 139 4881 2636}
  \quad\textbullet\quad
  Home Page: \href{https://bryceyin13.github.io/}{\textit{https://bryceyin13.github.io/}}
\end{center}

\vspace{-0.15\baselineskip}

\begin{cvsection}{Education}
\educationentry{École polytechnique fédérale de Lausanne (EPFL)}
  {Lausanne, Switzerland}
  {M.Sc. in Digital Humanities}
  {09/2025 - Present}
\begin{itemize}
  \item Research Focus: Computer Graphics and Computer Vision
  % \item Key Coursework: Advanced Computer Graphics, Applied Data Analysis
  \item \textbf{GPA:} 5.56 /6.00
\end{itemize}

\vspace{0.56\baselineskip}
\educationentry{Wuhan University}
  {Wuhan, China}
  {B.Eng. in Computer Science and Technology}
  {09/2021 - 06/2025}
\begin{itemize}
  \item \textbf{GPA:} 3.93 /4.00 (top 4\%)
\end{itemize}
\end{cvsection}

\begin{cvsection}{Publications}
\textbf{MATStruct: High-Quality Medial Mesh Computation via Structure-aware Variational Optimization}

Ningna Wang, Rui Xu, \textbf{\textit{Yibo Yin}}, Zichun Zhong, Taku Komura, Wenping Wang, Xiaohu Guo

\textit{ACM SIGGRAPH Asia 2025 (Conference Track)}
\end{cvsection}

\begin{cvsection}{Research Experience}
  \researchentry{Image and Visual Representation Lab, {\small EPFL}}{06/2026 - Present}
\textit{Research Assistant | {\small PI: Prof. Sabine Süsstrunk \externalicon{https://people.epfl.ch/sabine.susstrunk?lang=en}; Supervisor: Zhuoqian Yang, 
Doctoral Assistant \externalicon{https://inrainbws.github.io/}}}
\begin{itemize}
  \item Investigating weight space generation for implicit neural representations, with the goal of learning modality-agnostic priors over neural fields across multiple data modalities.
\end{itemize}

\vspace{0.58\baselineskip}
\researchentry{Realistic Graphics Lab, {\small EPFL}}{03/2026 - 06/2026}
\textit{Research Assistant | {\small Advisor: Prof. Wenzel Jakob \externalicon{https://rgl.epfl.ch/people/wjakob}}}
\begin{itemize}
  \item Developed a differentiable Gaussian Process Implicit Surface (GPIS) sampling module in Dr.Jit to support inverse-rendering with implicit geometry.
  \item Developed a hybrid surface-volume rendering approach for Many-Worlds inverse rendering.
\end{itemize}

\vspace{0.58\baselineskip}
\researchentry{Computer Graphics Lab, {\small The University of Texas at Dallas}}{04/2024 - 11/2025}
\textit{Research Assistant (Remote) | {\small Advisors: Prof. Xiaohu Guo \externalicon{https://personal.utdallas.edu/~xguo/} and Dr. Ningna Wang \externalicon{https://ningnawang.github.io/}}}
\begin{itemize}
  \item Adapted Loop subdivision surface evaluation and sharp feature fitting algorithms to structured medial mesh.
  \item Designed a simplification algorithm for structured medial mesh to preserve geometric structure under subdivision.
  \item Proposed a subdivision surface fitting method for structured medial mesh, targeting sharp-feature preservation during fitting.
\end{itemize}

\vspace{0.58\baselineskip}
\researchentry{Graphics and Vision Lab, {\small Wuhan University}}{05/2023 - 01/2024}
\textit{Research Assistant | {\small Advisor: Prof. Chunxia Xiao \externalicon{http://graphvision.whu.edu.cn/}}}
\begin{itemize}
  \item Built an RGB-ToF/depth data acquisition pipeline and collected 10,000+ paired image-sensor samples.
  \item Proposed a novel monocular depth estimation method using RGB images and auxiliary ToF/depth signals.
\end{itemize}
\end{cvsection}

\begin{cvsection}{Selected Projects}
\projectentry{Physically Based Ray Tracer}{2026}
\textit{CS-440 Advanced Computer Graphics Rendering Competition, EPFL}
\begin{itemize}
  \item Extended the Nori framework with heterogeneous volumetric path tracing, delta and ratio tracking, Henyey-Greenstein phase functions, and MIS over emitter, phase, and BSDF sampling.
  \item Implemented HDR environment map lighting, thin-lens depth of field, motion blur, principled BSDF, and GGX rough dielectric materials.
\end{itemize}

% \vspace{0.58\baselineskip}
% \researchentry{Waterloo Computer Graphics Lab, {\small University of Waterloo}}{06/2024 - 08/2024}
% \textit{Research Project (Remote) | {\small Advisor: Prof. Toshiya Hachisuka \externalicon{https://cs.uwaterloo.ca/~thachisu/}}}
% \begin{itemize}
%   \item Reproduced the algorithm in the SIGGRAPH paper \textit{A Practical Walk-on-Boundary Method for Boundary Value Problems} for solving the boundary value problem of Laplace's equation with Dirichlet boundaries.
% \end{itemize}

% \vspace{0.58\baselineskip}
% \projectentry{Photon Mapping Renderer}{2023}
% \textit{Personal Project}
% \begin{itemize}
%   \item Implemented photon mapping with final gathering based on \textit{Global Illumination using Photon Maps}, including photon emission, kd-tree nearest-neighbor queries, radiance estimation, and indirect illumination rendering.
% \end{itemize}

\end{cvsection}

% \newpage

\begin{cvsection}{Awards}
\awardentry{\textbf{Outstanding Student} (10\% school-wide), Wuhan University}{2022, 2023, 2024}
\awardentry{\textbf{Second Class Scholarship} (10\% school-wide), Wuhan University}{2022}
\awardentry{\textbf{Third Class Scholarship} (15\% school-wide), Wuhan University}{2023, 2024}
\awardentry{\textbf{Lei Jun (Xiaomi) Computer Innovation and Development Fund}, Wuhan University}{2024}
\end{cvsection}

\begin{cvsection}{Skills}
\begin{itemize}
  \item \textbf{Languages:} Mandarin (Native), English (TOEFL iBT 102)
  \item \textbf{Programming Languages:} C++, C, Python, C\#, GLSL, SQL, Verilog HDL, Java, \LaTeX
  \item \textbf{Libraries / Frameworks / Tools:} PyTorch, Dr.Jit, Mitsuba, PBRT-v3, Blender, Eigen, CGAL, libigl, Git, CMake
\end{itemize}
\end{cvsection}

\end{document}
